ANIMA is an experimental developmental artificial nervous system: a
deterministically seeded spiking substrate in which only the boundaries
--- sensory input, output, environment, curriculum, resource limits,
and local plasticity laws --- are specified, and nothing else. Phase I
asked whether useful, stable, adaptive internal organization can arise
from such a minimal specification, and progressed through a series of
frozen, identity-gated experiments to locate the architectural limit of
persistent learning. The record demonstrates, in turn: sensory-specific
synaptic organization emerges from spike-timing-dependent plasticity
under purely local rules; transient spike-based representations exist
during stimulation; scalar intrinsic persistence is information-poor
and cannot serve as sensory memory; the synaptic weight structure is
the substrate's persistent store, but it suffers order-dependent
interference --- blocked sequential experience overwrites, alternated
experience superposes; bounded consolidation protects the store from
runaway and preserves alternating coexistence; and, once stability was
established, the limiting factor moved successively to allocation, to
first-exposure allocation, and finally to the recruitment of postsynaptic
activity for genuinely novel late-arriving patterns, which Phase I
leaves unresolved. Throughout, the project enforced pre-registered
protocols, byte-identical flag-off control gates, deterministic replay,
and preservation of negative results, so that each negative narrows the
space of possible architectures rather than being discarded. The paper
reports what is demonstrated, what is supported, what is unmeasured,
and what remains hypothesis, with every quantitative claim traceable to
a committed repository record.